\chapter{1983年全国卷（理）}

\section{单选题}

\begin{problem}
两条异面直线，指的是（\quad）

\choices
    {在空间内不相交的两条直线}
    {分别位于两个不同平面内的两条直线}
    {某一平面内的一条直线和这个平面外的一条直线}
    {不在同一平面内的两条直线}

\begin{answer}
D
\end{answer}

\begin{solution}
异面直线的定义是既不相交又不在同一平面内的两条直线. 两条相交直线必在同一平面内，所以“不在同一平面内”已经同时排除了相交情形. 选项 D 与定义等价，而选项 A 只说明不相交，选项 B、C 都没有排除两直线共面的情形，因此正确选项为 D.
\end{solution}
\end{problem}

\begin{problem}
方程 \(x^{2} - y^{2} = 0\) 表示的图形是（\quad）

\choices
    {两条相交直线}
    {两条平行直线}
    {两条重合直线}
    {一个点}

\begin{answer}
A
\end{answer}

\begin{solution}
由平方差公式，\(x^{2}-y^{2}=(x-y)(x+y)=0\)，所以图形由直线 \(y=x\) 与直线 \(y=-x\) 组成. 这两条直线相交于原点，故选 A.
\end{solution}
\end{problem}

\begin{problem}
三个数 \(a\)，\(b\)，\(c\) 不全为零的充要条件是（\quad）

\choices
    {\(a\)，\(b\)，\(c\) 都不是零}
    {\(a\)，\(b\)，\(c\) 中最多有一个是零}
    {\(a\)，\(b\)，\(c\) 中只有一个是零}
    {\(a\)，\(b\)，\(c\) 中至少有一个不是零}

\begin{answer}
D
\end{answer}

\begin{solution}
“\(a\)，\(b\)，\(c\) 不全为零”就是“并非 \(a=b=c=0\)”. 由“全为零”的否定可知，这等价于 \(a\)，\(b\)，\(c\) 中至少有一个不是零，故正确选项为 D；其余选项都把条件误加强为零的个数受到更严格限制.
\end{solution}
\end{problem}

\begin{problem}
设 \(\alpha = \frac{4\pi}{3}\)，则 \(\arccos\left(\cos\alpha\right)\) 的值是（\quad）

\choices
    {\(\frac{4\pi}{3}\)}
    {\(-\frac{2\pi}{3}\)}
    {\(\frac{2\pi}{3}\)}
    {\(\frac{\pi}{3}\)}

\begin{answer}
C
\end{answer}

\begin{solution}
反余弦函数的值域为 \([0,\pi]\). 因为 \(\cos\frac{4\pi}{3}=\cos\frac{2\pi}{3}=-\frac{1}{2}\)，所以 \(\arccos\left(\cos\frac{4\pi}{3}\right)=\frac{2\pi}{3}\)，故选 C.
\end{solution}
\end{problem}

\begin{problem}
\(0.3^{2}\)，\(\log_{2}0.3\)，\(2^{0.3}\) 这三个数之间的大小顺序是（\quad）

\choices
    {\(0.3^{2} < 2^{0.3} < \log_{2}0.3\)}
    {\(0.3^{2} < \log_{2}0.3 < 2^{0.3}\)}
    {\(\log_{2}0.3 < 0.3^{2} < 2^{0.3}\)}
    {\(\log_{2}0.3 < 2^{0.3} < 0.3^{2}\)}

\begin{answer}
C
\end{answer}

\begin{solution}
因为 \(0<0.3<1\)，所以 \(\log_{2}0.3<0\). 又 \(0.3^{2}=0.09>0\)，而 \(2^{0.3}>1\). 因此 \(\log_{2}0.3<0.3^{2}<2^{0.3}\)，故选 C.
\end{solution}
\end{problem}

\section{解答题}

\begin{problem}
\begin{enumerate}
\item 在同一平面直角坐标系内，分别画出两个方程 \(y = -\sqrt{x}\)，\(x = \sqrt{-y}\) 的图形，并写出它们交点的坐标.
\item 在极坐标系内，方程 \(\rho = 5\cos \theta\) 表示什么曲线？画出它的图形.
\end{enumerate}

\begin{solution}
\begin{enumerate}
\item 对于 \(y=-\sqrt{x}\)，有 \(x\geqslant0\)，\(y\leqslant0\)，平方得 \(x=y^{2}\)，所以它是抛物线 \(x=y^{2}\) 的下半支. 对于 \(x=\sqrt{-y}\)，有 \(x\geqslant0\)，\(y\leqslant0\)，平方得 \(y=-x^{2}\)，所以它是抛物线 \(y=-x^{2}\) 的右半支.

令两式同时成立，则 \(y=-\sqrt{x}\)，\(y=-x^{2}\). 由 \(\sqrt{x}=x^{2}\) 且 \(x\geqslant0\)，令 \(u=\sqrt{x}\geqslant0\)，则 \(u=u^{4}\)，即 \(u(u^{3}-1)=0\)，所以 \(u=0\) 或 \(u=1\)，从而 \(x=0\) 或 \(x=1\)，相应地 \(y=0\) 或 \(y=-1\). 因此交点为 \((0,0)\) 和 \((1,-1)\).
\item 将极坐标关系 \(x=\rho\cos\theta\)，\(y=\rho\sin\theta\) 代入，并乘以 \(\rho\)，得
\[
x^{2}+y^{2}=\rho^{2}=5\rho\cos\theta=5x
\]
所以
\[
\left(x-\frac{5}{2}\right)^{2}+y^{2}=\frac{25}{4}
\]
该曲线是圆心为 \(\left(\frac{5}{2},0\right)\)、半径为 \(\frac{5}{2}\) 的圆，经过 \((0,0)\) 和 \((5,0)\)，图形据此作出.
\end{enumerate}
\end{solution}
\end{problem}

\begin{problem}
\begin{enumerate}
\item 已知 \(y = \e^{-x}\sin 2x\)，求微分 \(\mathrm{d}y\).
\item 一个小组共有 \(10\) 名同学，其中 \(4\) 名是女同学，\(6\) 名是男同学. 要从小组内选出 \(3\) 名代表，其中至少有 \(1\) 名女同学，求一共有多少种选法.
\end{enumerate}

\begin{solution}
\begin{enumerate}
\item 由乘积求导法则，
\[
\frac{\mathrm{d}y}{\mathrm{d}x}=\left(\e^{-x}\right)'\sin 2x+\e^{-x}(\sin 2x)'=-\e^{-x}\sin 2x+2\e^{-x}\cos 2x
\]
因此 \(\mathrm{d}y=\e^{-x}\left(2\cos 2x-\sin 2x\right)\mathrm{d}x\).
\item 按女同学人数分类. 恰有 \(1\) 名女同学时有 \(\mathrm C_{4}^{1}\mathrm C_{6}^{2}=4\times15=60\) 种；恰有 \(2\) 名女同学时有 \(\mathrm C_{4}^{2}\mathrm C_{6}^{1}=6\times6=36\) 种；\(3\) 名均为女同学时有 \(\mathrm C_{4}^{3}=4\) 种. 所以共有 \(60+36+4=100\) 种选法.
\end{enumerate}
\end{solution}
\end{problem}

\begin{problem}
计算行列式（要求结果最简）：\(\begin{vmatrix}\sin \alpha & \cos \left(\alpha + \varphi\right) & \cos \alpha \\ \cos \beta & \sin \left(\beta - \varphi\right) & \sin \beta \\ \sin \varphi & \cos 2 \varphi & \cos \varphi\end{vmatrix}\).

\begin{solution}
设行列式为 \(\Delta\). 记三列分别为 \(C_{1}\)，\(C_{2}\)，\(C_{3}\). 由三角函数加法公式，
\[
\cos(\alpha+\varphi)=\cos\varphi\cos\alpha-\sin\varphi\sin\alpha
\]
\(\sin(\beta-\varphi)=\cos\varphi\sin\beta-\sin\varphi\cos\beta\)，\(\cos2\varphi=\cos^{2}\varphi-\sin^{2}\varphi\).
故 \(C_{2}=\cos\varphi\,C_{3}-\sin\varphi\,C_{1}\)，三列线性相关，从而 \(\Delta=0\).
\end{solution}
\end{problem}

\begin{problem}
\begin{enumerate}
\item 证明：对于任意实数 \(t\)，复数 \(z = \sqrt{|\cos t|} + \sqrt{|\sin t|}\i\) 的模 \(r = |z|\) 适合 \(r \leqslant \sqrt[4]{2}\).
\item 当实数 \(t\) 取什么值时，复数 \(z = \sqrt{|\cos t|} + \sqrt{|\sin t|}\i\) 的幅角主值 \(\theta\) 适合 \(0 \leqslant \theta \leqslant \frac{\pi}{4}\)？
\end{enumerate}

\begin{solution}
\begin{enumerate}
\item 由复数模的定义，
\[
r^{2}=|z|^{2}=|\cos t|+|\sin t|
\]
对非负数 \(u=|\cos t|\)，\(v=|\sin t|\)，由 \((u-v)^{2}\geqslant0\) 得 \(u+v\leqslant\sqrt{2(u^{2}+v^{2})}\). 因为 \(u^{2}+v^{2}=\cos^{2}t+\sin^{2}t=1\)，所以 \(r^{2}\leqslant\sqrt{2}\)，即 \(r\leqslant\sqrt[4]{2}\). 当且仅当 \(|\cos t|=|\sin t|\) 时取等号.
\item \(z\) 的实部和虚部均非负，且不同时为零，因此其幅角主值满足 \(0\leqslant\theta\leqslant\frac{\pi}{2}\). 在此范围内，
\[
0\leqslant\theta\leqslant\frac{\pi}{4}\iff \sqrt{|\sin t|}\leqslant\sqrt{|\cos t|}\iff |\sin t|\leqslant|\cos t|
\]
若 \(\cos t=0\)，则 \(|\sin t|=1>|\cos t|=0\)，不满足上述不等式；若 \(\cos t\neq0\)，则
\[
|\sin t|\leqslant|\cos t|\iff |\tan t|\leqslant1
\]
在每个区间 \(k\pi-\frac{\pi}{2}<t<k\pi+\frac{\pi}{2}\) 内，\(\tan t\) 严格递增，因此 \(|\tan t|\leqslant1\) 等价于
\(k\pi-\frac{\pi}{4}\leqslant t\leqslant k\pi+\frac{\pi}{4}\)，\(k\in\Z\).
\end{enumerate}
\end{solution}
\end{problem}

\begin{problem}
如图，在三棱锥 \(S - ABC\) 中，\(S\) 在底面上的射影 \(N\) 位于底面的高 \(CD\) 上；\(M\) 是侧棱 \(SC\) 上的一点，使截面 \(MAB\) 与底面所成的角等于 \(\angle NSC\)，求证：\(SC\) 垂直于截面 \(MAB\).

\bitmapfigure[width=4.9cm]{img/1983/national_science/q9_fig1.png}

\begin{solution}
因为 \(SN\) 是 \(S\) 在底面上的射影，所以 \(SN\perp AB\)；又因为 \(CD\) 是底面三角形的高，所以 \(CD\perp AB\). 在平面 \(SCD\) 内过 \(D\) 作直线 \(\ell\parallel SN\)，则 \(\ell\perp AB\)，且 \(CD\perp AB\). 由于 \(\ell\) 与 \(CD\) 是平面 \(SCD\) 内相交的两条直线，故 \(AB\perp\left(SCD\right)\)，从而 \(AB\perp SC\).

平面 \(MAB\) 与底面相交于 \(AB\). 因为 \(D\in AB\)，且 \(M\)，\(D\in\left(SCD\right)\)，所以 \(MD\) 同时位于平面 \(MAB\) 和平面 \(SCD\) 内. 又 \(MD\perp AB\)、\(CD\perp AB\)，故截面与底面所成的角为 \(\angle MDC\). 由题设，\(\angle MDC=\angle NSC\). 在平面 \(SCD\) 内，\(SN\perp CD\)，因此
\[
\angle NSC+\angle SCD=\frac{\pi}{2}
\]
因为 \(M\in SC\)、\(N\in CD\)，在三角形 \(MDC\) 中有 \(\angle MCD=\angle SCD\). 所以 \(\angle MDC+\angle MCD=\frac{\pi}{2}\)，从而 \(\angle DMC=\frac{\pi}{2}\)，即 \(MD\perp MC\)，也就是 \(MD\perp SC\). 直线 \(AB\) 与 \(MD\) 在平面 \(MAB\) 内相交，且 \(SC\) 同时垂直于它们，所以 \(SC\) 垂直于平面 \(MAB\).
\end{solution}
\end{problem}

\begin{problem}
如图，已知椭圆长轴 \(\left|A_{1}A_{2}\right| = 6\)，焦距 \(\left|F_{1}F_{2}\right| = 4\sqrt{2}\)，过椭圆焦点 \(F_{1}\) 作一直线，交椭圆于两点 \(M\)，\(N\). 设 \(\angle F_{2}F_{1}M = \alpha\)（\(0 \leqslant \alpha < \pi\)），当 \(\alpha\) 取什么值时，\(|MN|\) 等于椭圆短轴的长？

\bitmapfigure[width=6.2cm]{img/1983/national_science/q10_fig1.png}

\begin{solution}
设椭圆中心为原点，长轴在 \(x\) 轴上. 由题意，半长轴 \(a=3\)，半焦距 \(c=2\sqrt{2}\)，所以半短轴
\[
b=\sqrt{a^{2}-c^{2}}=\sqrt{9-8}=1
\]
椭圆方程为 \(\frac{x^{2}}{9}+y^{2}=1\)，且 \(F_{1}=(-2\sqrt{2},0)\). 过 \(F_{1}\) 的直线方向向量可取 \((\cos\alpha,\sin\alpha)\)，其参数方程为
\(x=-2\sqrt{2}+t\cos\alpha\)，\(y=t\sin\alpha\).
代入椭圆方程，得到
\[
\frac{1+8\sin^{2}\alpha}{9}t^{2}-\frac{4\sqrt{2}\cos\alpha}{9}t-\frac{1}{9}=0
\]
设两根为 \(t_{1}\)，\(t_{2}\). 因为参数方向向量是单位向量，弦长 \(|MN|=|t_{1}-t_{2}|\). 该二次方程的判别式为
\[
\left(-\frac{4\sqrt{2}\cos\alpha}{9}\right)^{2}
-4\cdot\frac{1+8\sin^{2}\alpha}{9}\cdot\left(-\frac{1}{9}\right)
=\frac{32\cos^{2}\alpha+4+32\sin^{2}\alpha}{81}
=\frac{4}{9}
\]
所以
\[
|MN|=\frac{\sqrt{4/9}}{(1+8\sin^{2}\alpha)/9}=\frac{6}{1+8\sin^{2}\alpha}
\]
椭圆短轴长为 \(2b=2\)，故
\[
\frac{6}{1+8\sin^{2}\alpha}=2\iff \sin^{2}\alpha=\frac14
\]
又 \(0\leqslant\alpha<\pi\)，所以 \(\alpha=\frac{\pi}{6}\) 或 \(\alpha=\frac{5\pi}{6}\).
\end{solution}
\end{problem}

\begin{problem}
已知数列 \(\{a_{n}\}\) 的首项 \(a_{1} = b\)（\(b \neq 0\)），它的前 \(n\) 项的和 \(S_{n} = a_{1} + a_{2} + \cdots + a_{n}\)（\(n \geqslant 1\)），并且 \(S_{1}\)，\(S_{2}\)，\(S_{n}\)，\(\cdots\) 是一个等比数列，其公比为 \(p\)（\(p \neq 0\)，\(|p| < 1\)）.

\begin{enumerate}
\item 证明：\(a_{2}\)，\(a_{3}\)，\(\cdots\)，\(a_{n}\)，\(\cdots\)（即 \(\{a_{n}\}\) 从第二项起）是一个等比数列；
\item 设 \(W_{n}=a_{1}S_{1}+a_{2}S_{2}+a_{3}S_{3}+\cdots+a_{n}S_{n}\)（\(n\geqslant1\)），求 \(\displaystyle\lim_{n\to\infty}W_{n}\)（用 \(b\)，\(p\) 表示）.
\end{enumerate}

\begin{solution}
\begin{enumerate}
\item 因为 \(S_{1}=a_{1}=b\)，且 \(S_{1}\)，\(S_{2}\)，\(\ldots\) 是公比为 \(p\) 的等比数列，所以 \(S_{n}=bp^{n-1}\). 对 \(n\geqslant2\)，
\[
a_{n}=S_{n}-S_{n-1}=bp^{n-2}(p-1)
\]
由于 \(b\neq0\)、\(p\neq0\) 且 \(|p|<1\)，有 \(p\neq1\)，所以 \(a_{n+1}/a_{n}=p\). 因而 \(a_{2}\)，\(a_{3}\)，\(\ldots\) 是公比为 \(p\) 的等比数列.
\item \(a_{1}S_{1}=b^{2}\). 对 \(k\geqslant2\)，
\[
a_{k}S_{k}=bp^{k-2}(p-1)\cdot bp^{k-1}=b^{2}(p-1)p^{2k-3}
\]
当 \(n\geqslant2\) 时，
\[
W_{n}=b^{2}+b^{2}(p-1)p\sum_{j=0}^{n-2}p^{2j}
\]
因为 \(|p|<1\)，所以
\[
\lim_{n\to\infty}W_{n}=b^{2}+\frac{b^{2}(p-1)p}{1-p^{2}}=b^{2}-\frac{b^{2}p}{1+p}=\frac{b^{2}}{1+p}
\]
\end{enumerate}
\end{solution}
\end{problem}

\begin{problem}
\begin{enumerate}
\item 已知 \(a\)，\(b\) 为实数，并且 \(\e < a < b\)，其中 \(\e\) 是自然对数的底数，证明：\(a^b > b^a\)；
\item 如果正实数 \(a\)，\(b\) 满足 \(a^b = b^a\)，且 \(a < 1\)，证明：\(a = b\).
\end{enumerate}

\begin{solution}
\begin{enumerate}
\item 令 \(f(x)=\frac{\ln x}{x}\)，则
\[
f'(x)=\frac{1-\ln x}{x^{2}}
\]
当 \(x>\e\) 时，\(f'(x)<0\)，故 \(f\) 在 \((\e,+\infty)\) 上单调递减. 由 \(\e<a<b\)，得 \(f(a)>f(b)\)，于是
\[
\ln\frac{a^{b}}{b^{a}}=b\ln a-a\ln b=ab\left(\frac{\ln a}{a}-\frac{\ln b}{b}\right)>0
\]
因此 \(a^{b}>b^{a}\).
\item 由 \(a^{b}=b^{a}\) 及 \(a\)，\(b>0\)，取自然对数得 \(b\ln a=a\ln b\). 因为 \(0<a<1\)，所以 \(\ln a<0\)，从而等式右边也小于零，故 \(\ln b<0\)，即 \(0<b<1\). 函数 \(f(x)=\frac{\ln x}{x}\) 在 \((0,1)\) 上满足
\[
f'(x)=\frac{1-\ln x}{x^{2}}>0
\]
所以 \(f\) 在 \((0,1)\) 上严格递增. 等式 \(b\ln a=a\ln b\) 等价于 \(f(a)=f(b)\)，因而 \(a=b\).
\end{enumerate}
\end{solution}
\end{problem}
